\section{Edition 05: entry, fixed-type attainment and rigidity}
\label{sec:ed05}
The open problems describe the geometric targets. The checked propositions below give their current formal prerequisites. Fix a closed connected oriented
smooth homotopy four-sphere $M$. Put $T=4\pi^2$,
$e=E_+=E_-$ on closed metrics, and $I=I^+(M)$.
The equality of the two closed integrals does not identify their restrictions
to a subdomain. All energy and volume measurements use their supplying metric.

\begin{openproblem}[Positive entry on a specified smooth type]
\label{prop:ed05-entry}
On a specified $M$ construct a positive Yamabe class with $e<2T$.
One target is a smooth representative with $R>0$ and
$|\operatorname{Ric}_0|^2<R^2/12$ everywhere. For a marked handle
boundary establish the smooth identification and the estimates on the entire
closed representative. Record which smooth types the construction covers.
\end{openproblem}

\begin{openproblem}[Budget on the original regularized trajectory]
\label{prop:ed05-bour-energy}
\uses{lem:ed03-positive-sequence,prop:ed47-budget,prop:ed47-motion}
Given entry, choose $0<\alpha_j<1$ decreasing to zero and unit-volume
positive Yamabe metrics with $0\le e_{j,0}-I\le\alpha_j^2$. Write
$D=\int|\operatorname{Ric}_0|^2$ and $S=\int R^2$. Along the original
negative twice-gradient flow of
$\mathcal F^\alpha=2(1-\alpha)e+\alpha D/2$,
prove that for all sufficiently large $j$ and every smooth time,
$Y\ge y>0$, $I\le e(t)\le 2T-\eta$ and
$D(t)\le4I+4(e_{j,0}-I)/\alpha_j\le d_*<32\pi^2$,
with $y,\eta,d_*$ independent of $j$ and time.
\end{openproblem}

\begin{openproblem}[The two outputs of one fixed-parameter flow]
\label{prop:ed05-bour-alternative}
\uses{prop:ed05-bour-energy}
At a fixed sufficiently small positive $\alpha_j$, instantiate Bour's
strictly pinched alternative on that original maximal solution. Return either
a closed $\mathcal F^{\alpha_j}$-critical metric via whole smooth comparison
or a complete nonflat scalar-flat Bach-flat pointed limit with the
physical bounds and coherent embeddings on compact domains. Account for
the short-time gauge and derivative estimates at fixed parameter.
\end{openproblem}

\begin{openproblem}[Smooth identification of a pointed limit]
\label{prop:ed05-bour-pointed}
\uses{prop:ed47-mark,prop:ed47-ends}
Take a smooth pointed limit of rescaled original metrics with the uniform
positive Yamabe, subcritical chiral and strict Ricci bounds of the preceding
budget, finite full curvature, and $R=B=0$ on the complete nonflat limit.
Use one oriented embedding on a neighborhood of the connected truncated
core and every end buffer. Exclude nontrivial end groups before using
$D_X=2(E_+(X)+E_-(X))+32\pi^2(k-1)$ to deduce $k=1$.
Construct the smooth conformal compactification $N$ and its positive
class. Identify $N$ with $M$ by forward and inverse smooth maps
agreeing on the actual cap collar. Prove $I\le e_X\le\liminf_i e_i$.
\end{openproblem}

\begin{openproblem}[Uniform estimates for the mixed critical sequence]
\label{prop:ed05-bour-stationary}
\uses{prop:ed05-bour-alternative,prop:ed05-bour-pointed,prop:ed47-balls,prop:ed47-mixed}
Return one unit-volume constant-scalar metric $q_j$ on $M$ at each late
index. Its parameter is $p_j=\alpha_j$ for a closed regularized return
and $p_j=0$ for an exact pointed return in Yamabe gauge.
Recover the scalar, physical Ricci and energy bounds in the gauges used.
From the constant-scalar trace-free Ricci equation with principal term
$\tfrac12\nabla^*\nabla\operatorname{Ric}_0$, prove noncollapse,
harmonic-coordinate regularity and estimates of every order uniform
as $p_j\to0$, also on curvature-normalized balls in the unbounded branch.
\end{openproblem}

\begin{openproblem}[Closed return at bounded curvature]
\label{prop:ed05-bour-closed}
\uses{prop:ed05-bour-stationary,prop:ed47-return}
If the unscaled curvature of the selected critical sequence is bounded
on a subsequence, obtain smooth convergence with whole oriented
comparison diffeomorphisms. Pass the positive Yamabe floor, the
vanishing Bach limit and $e(q_j)\to I$ to the limit and pull back to $M$.
\end{openproblem}

\begin{openproblem}[An attained positive Bach-flat class on the entered type]
\label{thm:ed05-bour-attainment}
\uses{prop:ed05-bour-closed,prop:ed05-bour-pointed,prop:ed05-bour-stationary}
From one positive subcritical entry on $M$ construct a smooth positive
Bach-flat class on $M$ of energy exactly $I$. In the unbounded-curvature
branch of the mixed critical sequence, use its own nonflat normalized
pointed limit, the marked return, and $I\le e_X\le\liminf e(q_j)=I$.
\end{openproblem}

\begin{openproblem}[Uniform infimum gap on an entered type]
\label{prop:ed05-uniform-gap}
\uses{prop:master-recognition}
Realize the master's uniform recognition gap as a separate formal
dependency. There is an $\varepsilon_*>0$, independent of the
oriented smooth homotopy four-sphere $M$, such that a nonempty
positive-class energy set with $I^+(M)\le T+\varepsilon_*$ implies
that $M$ is diffeomorphic to the standard sphere. Trace the
positive-time smoothing, compactness and smooth type hypotheses.
\end{openproblem}

\begin{openproblem}[Rigidity of an attained minimizer]
\label{prop:ed05-rigidity}
\uses{prop:ed47-variation}
Assume a smooth positive class attains $I$ in the band above the
recognition gap and below $2T$. For a specified candidate prove smooth
recognition of $M$ or produce a strictly lower positive competitor on
$M$. In a Hessian test exhibit an actual admissible negative direction
modulo gauge. Record the surviving Einstein and non-Einstein cases.
\end{openproblem}

\begin{openproblem}[Coverage of the smooth types]
\label{thm:ed05-coverage}
\uses{prop:ed05-entry,thm:ed05-bour-attainment,prop:ed05-rigidity,prop:ed05-uniform-gap,prop:master-recognition,prop:master-immersion}
For every oriented smooth homotopy four-sphere give either an applicable
direct recognition certificate or an entry witness and a rigidity
argument covering the minimizing class returned on that same type.
For a handle argument preserve the actual markings and establish
exhaustiveness of its allowed presentations and successors.
\end{openproblem}

\subsection{Checked results for the original and returned metrics}

\begin{proposition}[Yamabe starts at prescribed excess]
\label{prop:ed47-start}
\lean{SmoothPoincare.BourYamabeNormalization, SmoothPoincare.exists_bourYamabeStarts}
\leanok
Given one measured positive entry below $8\pi^2$ and a conformal Yamabe
normalization preserving the measured energy, every positive sequence
$\alpha_j$ admits unit-volume positive Yamabe starts on this carrier.
Their excess lies in $[0,\alpha_j^2]$, their scalar-square integral is
$Y_{j,0}^2$, and they satisfy a common strict chiral margin.
\end{proposition}
\begin{proof}
\leanok
Use the entry to inhabit the positive energy set. Select a member below the
smaller of $I+\alpha_j^2$ and the midpoint between the entry energy and
$8\pi^2$. Normalize the selected metric in its own class.
\end{proof}

\begin{proposition}[Physical budget at every original time]
\label{prop:ed47-budget}
\uses{prop:ed47-start}
\lean{SmoothPoincare.Bour.bourChargeOn_full_eq, SmoothPoincare.Bour.bourEuler_eq_charges, SmoothPoincare.Bour.bourEnergy_eq_chiral_scalar, SmoothPoincare.Bour.bourEnergy_conformal_minimum, SmoothPoincare.Bour.bourRicci_eq_integral, SmoothPoincare.Bour.bourFullCurvature_eq_euler_add_ricci, SmoothPoincare.Bour.bourOriginalFlow_budget, SmoothPoincare.Bour.eventually_bourOriginalFlow_budget}
\leanok
Supply the closed Euler value $16\pi^2$, the signed metric measurements,
the unit-volume Yamabe representatives, the round initial upper bound,
a continuous sign detector and the original functional sublevel.
Put $\delta=8\pi^2-I>0$ and $\epsilon_j=e_{j,0}-I\le\alpha_j^2$.
When $0<\alpha_j<1$ tends to zero, one cutoff in $j$ gives, at every
original time,
\[
 Y_t\ge\sqrt{24\delta},\qquad I\le E_\pm(g_t)\le8\pi^2-\delta/2,
 \qquad S_t\le Y_{j,0}^2+48\alpha_j .
\]
The physical two-tensor norm satisfies
\[
 D_t\le4I+4\epsilon_j/\alpha_j+(Y_{j,0}^2-384\pi^2)/12
 \le32\pi^2-2\delta.
\]
The curvature-operator integral is $16\pi^2+D_t$.
Both total chiral energies converge uniformly over the original time domains to $I$.
\end{proposition}
\begin{proof}
\leanok
Integrate the local four-slot curvature decomposition before using closed
signature. The forms $2(1-\alpha)e+\alpha D/2$ and
$2e+\alpha S/24-16\alpha\pi^2$ agree. A strictly positive pinching
margin bounds the Yamabe square away from zero. Connectedness propagates
the initial sign. Every slice is now admissible in the same infimum, which
gives the displayed inequalities by functional decrease.
\end{proof}

\begin{proposition}[Dissipation and the squared parabolic time scale]
\label{prop:ed47-motion}
\lean{SmoothPoincare.Bour.original_gradient_sublevel, SmoothPoincare.Bour.metricVelocity_scale_time, SmoothPoincare.Bour.gradientSq_scale, SmoothPoincare.Bour.parabolic_dissipation_eq, SmoothPoincare.Bour.bourEnergy_scale, SmoothPoincare.Bour.volume_constant_of_scale_variation}
\leanok
For a supplied geometric first variation and actual velocity $-2G$ on
a regular slab, $\int_a^b\|G\|_2^2=(F_a-F_b)/2$ and $F_b\le F_a$.
Under $c g(t_0+s/c^2)$ the covariant velocity has weight $c^{-1}$ and
the spatial gradient integral has weight $c^{-2}$. The time-integrated
dissipation on corresponding slabs agrees. The closed energy is
constant under scaling at fixed Euler value. Volume is constant given
the regular volume variation, its velocity trace and scale-direction
first variation.
\end{proposition}
\begin{proof}
\leanok
Insert the metric derivative in the physical tensor pairing and integrate
the derivative of $F$. Four-volume scales by $c^2$. The time substitution
cancels the remaining power. The derivative of the functional in the
scaling direction vanishes, giving zero integrated velocity trace.
\end{proof}

\begin{proposition}[One comparison of the full truncation]
\label{prop:ed47-mark}
\lean{SmoothPoincare.bourWholeCoreComparison, SmoothPoincare.bourComparison_orientationAt, SmoothPoincare.bourPointedTraceFreeRicci_real_le_liminf, SmoothPoincare.bourPointedSignedWeyl_le_liminf, SmoothPoincare.bourGlueMarkedInverse}
\leanok
Given pointed exhaustion maps with the specified orientation, a compact
truncation containing the core and all buffers is eventually covered by
one open comparison and its inverse. Canonical two-jet metric convergence,
actual constant rescalings of compact sources and an eventual Ricci ceiling
give $D_X\le\liminf D_i$ in the original physical norm. Each Weyl sign
has its separate comparison under the signed frame and density hypotheses.
Two supplied partial diffeomorphisms on the core and its exterior, with
the prescribed agreement, image overlap, coverage and orientation at a
core point, glue to an oriented diffeomorphism with both inverse germs.
\end{proposition}
\begin{proof}
\leanok
Restrict the single late comparison. Exhaust the limit with compact sets
before bounding the full lower limit. The Ricci density has scale weight
$-2$ and the four-volume has weight $2$. Glue the two compact neighborhoods,
globalize using source and image coverage, and restrict the inverse.
\end{proof}

\begin{proposition}[End count and both capped exterior inequalities]
\label{prop:ed47-ends}
\lean{SmoothPoincare.bourOneEuclideanEnd, SmoothPoincare.bourSignedEndEnergy, SmoothPoincare.bourCappedExterior_both_lt}
\leanok
After triviality of all end groups and the spherical boundary Euler
identity have been supplied, $k\ge1$ and $D_X<32\pi^2$ imply $k=1$.
Equal signed totals on the compactified carrier then give $D_X=4e_X$.
For a specified cap the two signed exterior tests give strict capped
energies below $T$ under their cut, ball-mass and splice margins.
\end{proposition}
\begin{proof}
\leanok
The boundary identity is $D_X=2(E_++E_-)+32\pi^2(k-1)$. Nonnegative
signed charges exclude $k\ge2$. Subtract the signed interior masses and
add the measured splice on each complete exterior.
\end{proof}

\begin{proposition}[Uniform physical ball volumes for the CSC sequence]
\label{prop:ed47-balls}
\lean{SmoothPoincare.bourYamabe_smooth_test_scale, SmoothPoincare.bourBallVolume_lower_of_cutoff, SmoothPoincare.bourBallVolume_lower_of_yamabe_scale, SmoothPoincare.bourTraceFreeRicci_scale, SmoothPoincare.bourVolume_scale}
\leanok
On the connected closed four-manifold let $Y\ge y>0$, with constant
scalar $0\le R\le C$. For every $a\ge1$, center $x$ and $0<r\le1$,
\[
 \operatorname{Vol}_{a g} B_{a g}(x,r)\ge
 \left(\frac{y}{4(6L^2+C)}\right)^2 r^4,
\]
where $L$ is the fixed smooth cutoff slope. The constant works on the
original and curvature-normalized metrics at every index.
\end{proposition}
\begin{proof}
\leanok
Use the conformal-class inequality on nonnegative smooth tests. In the
scaled metric smooth a Lipschitz cutoff, bound its gradient by $L/r$,
then iterate the ball-volume inequality from the small-radius lower bound.
\end{proof}

\begin{proposition}[Initial Yamabe saturation in a second pointed return]
\label{prop:ed47-mixed}
\uses{prop:ed47-budget,prop:ed47-mark,prop:ed47-ends}
\lean{SmoothPoincare.bourFirstPointed_excess_and_deficit, SmoothPoincare.Bour.fixedReturn_budget, SmoothPoincare.Bour.mixedReturn_energy_and_parameter, SmoothPoincare.Bour.mixedReturn_pointed_saturation}
\leanok
Suppose the closed first returns give convergence of their physical
energies, and the pointed first returns give the marked Weyl and Ricci
comparisons and the scalar bound in their new Yamabe gauge. Positivity,
the closed Euler evaluations and original budgets supply $p_j\to0$ and
$E_\pm(q_j)\to I$. On a second marked pointed subsequence, let the
positive returned metric have $e_X\le\liminf E_+(q_j)$ and
$4e_X\le\liminf D(q_j)$. Then $e_X=I$ and, on that subsequence,
$Y_{j,0}^2\to384\pi^2$.
\end{proposition}
\begin{proof}
\leanok
Write $b_j=384\pi^2-Y_{j,0}^2$. A first pointed return gives
$b_j\le48\alpha_j$. A closed one gives
$b_j\le48\alpha_j+12(4I-D(q_j))$. The oriented positive return gives
$I\le e_X$. Use $4I\le\liminf D(q_j)$ to control the closed terms
even when the two first alternatives alternate.
\end{proof}

\begin{proposition}[Positive return on the original smooth carrier]
\label{prop:ed47-return}
\lean{SmoothPoincare.SignedCollarMetric.pullbackCross, SmoothPoincare.SignedCollarMetric.energy_tendsto_of_metricCPConvergence, SmoothPoincare.bourAttainsOfSignedPullback, SmoothPoincare.bourBachFlat_of_criticalLimit, SmoothPoincare.bourClosedReturn, SmoothPoincare.bourRoundReturn}
\leanok
A supplied oriented diffeomorphism pulls a positive measured metric of
energy at most the source infimum back to an attaining metric in both
chiralities. Bach-flat transport uses the specified covariant operator
and naturality input. On a supplied closed mixed return the energy,
conformal test and covariant Euler limits give the same conclusion. On
a fixed carrier, the two-jet energy comparison uses compatible frames,
weighted measurability and domination. The actual round pullback has both
energies zero.
\end{proposition}
\begin{proof}
\leanok
Transport the signed geometric witness and use membership in the positive
energy set for the opposite inequality. On the closed limit pass each
conformal test separately and take the limit of the supplied Euler identity.
\end{proof}

\begin{proposition}[Second variation at an attained geometric metric]
\label{prop:ed47-variation}
\uses{prop:ed47-return}
\lean{SmoothPoincare.bourAttainedVariation, SmoothPoincare.bourNegativeSecondVariationExcludesAttainment, SmoothPoincare.bourRoundAttainedVariation}
\leanok
Along a two-sided, jointly smooth positive metric variation with twice
continuously differentiable measured Weyl energy, an attaining metric
has zero first and nonnegative second derivative. Negative second
derivative in a supplied direction produces arbitrarily close positive
competitors of strictly lower energy. The signed round pullback satisfies
the necessary variation conditions.
\end{proposition}
\begin{proof}
\leanok
Positive-class membership gives a local minimum along the path.
Apply the second-derivative test to the real energy function.
\end{proof}
